\documentclass[10pt,a4paper]{amsart}
\usepackage[margin=19mm]{geometry}
\usepackage{amsmath,amssymb,mathtools}
\usepackage[scr=rsfs,cal=euler]{mathalfa}
\usepackage{xfrac}
\usepackage[unicode,hidelinks,bookmarks=false]{hyperref}
\newcommand{\ie}{i.e.\ }
\newcommand{\cf}{cf.\ }
\newcommand{\resp}{respectively\ }
\newcommand{\Subsection}{\subsection}
\providecommand{\nobreakdash}{\nobreak\hbox{-}}
\setlength{\emergencystretch}{2em}
\multlinegap0pt
\usepackage{float}
\usepackage{xcolor}
\usepackage{enumerate}
\usepackage{tikz}
\usepackage{amssymb}
\usepackage{dsfont}
\newtheorem{lemma}{Lemma}
\usepackage{cleveref}
\crefname{lemma}{Lemma}{Lemmas}
\newcommand{\C}{\mathbb{C}}
\newcommand{\R}{\mathbb{R}}
\newcommand{\sm}{\setminus}
\newcommand{\sse}{\subset}
\newcommand{\wt}{\widetilde}
\title{The Legendrian Hopf Link has exactly two Lagrangian fillings}
\author{Bryce Thomson}
\date{Source: arXiv:2506.15111v2 (CC BY 4.0)}
\begin{document}
\maketitle

\noindent\emph{Source and licence.} The Legendrian Hopf Link has exactly two Lagrangian fillings. Version arXiv:2506.15111v2 (CC BY 4.0), distributed by arXiv under the Creative Commons Attribution 4.0 International licence, http://creativecommons.org/licenses/by/4.0/, as recorded in the arXiv licence metadata for this exact version and shown on the abstract page https://arxiv.org/abs/2506.15111v2. Full licence notice: \texttt{licenses/arxiv-2506.15111-CC-BY-4.0.txt}.\par\smallskip

\noindent\textit{Editorial scope.} Counted unit: the article's lemma lem:horiziso (source0.tex lines 1029--1031). The article's complete original proof of it is delivered with it (source0.tex lines 1033--1060, one \texttt{proof} environment of 28 lines). No mathematics is written, repaired or re-derived: the statement and the proof are the article's own lines, in the article's own notation, and no retained line is substituted or re-flowed.  Retained uncounted as the source-local context the proof needs: lines 609--609; lines 702--703; lines 1011--1012.  One declared adaptation: exactly 1 ancillary strings inside the retained spans are omitted (image-inclusion commands and, where the source repeats a label, the repeated label definitions) and no other line differs from the source; the figure environments, with their placement, captions and labels, are kept, and the package records the omitted strings in fidelity receipts bound to the affected bodies.  No retained line contains a citation, so the wrapper carries no bibliography and no bibliography entry.  The retained lines contain the article's own diagram code, which the wrapper typesets with the standard tikz packages; no image file is delivered and the package declares no asset dependency.  Editorial formatting only, in addition: an AMS \texttt{amsart} preamble that restates the article's own declarations and macros used by the retained lines, namely the environments \texttt{lemma} on separate counters, together with the standard packages those lines need.  The retained environments print in this document as Lemma 1 in order of appearance, whereas the article prints the same items with the numbering of its own full document.  The complete native source of the article as distributed in the arXiv e-print for this exact version is bundled unchanged as \texttt{sources/180060/source0.tex}.
\subsection{Constructing Standard Lagrangian Fillings}\label{sec:std_fillings}

\begin{lemma}\label{lem:compends} The ends of $L$, $Ends(L)$, are compatible with the standard Lefschetz fibration $\tilde{f}:(B_1^4,\omega_0)\to \C$ where $\sigma_{Ends}:=\tilde{f}(Ends (L))$ are two disjoint properly embedded open curves, one near $\tilde{f}(\Lambda_1)=e^{i\pi/4}\in \C$ and the other near $\tilde{f}(\Lambda_2)=e^{i7\pi/4}\in \C$. Furthermore, $L$ intersects any of the fibers $f^{-1}(s)$ for $s\in \sigma_{\text{Ends}}$ precisely along the zero section.
\end{lemma}

\begin{lemma}\label{lem:hamisotozerosection} Suppose an embedded Lagrangian filling $L$ of $\Lambda_h$ is compatible with the standard Lefschetz fibration $\Tilde{f}:(B_1^4,\omega_0)\to \C$ over the curve $\sigma=\tilde{f}(L)\sse \C\sm \{0\}$. Then there exists a Hamiltonian isotopy so that $L$ intersects each fiber over $\sigma$ precisely along the zero section.
\end{lemma}

\begin{lemma}\label{lem:horiziso}
Suppose an embedded Lagrangian filling $L$ of $\Lambda_h$ is compatible with the standard Lefschetz fibration $\tilde{f}:(B^4_1,\omega_0)\to \C$ over the curve $\sigma=\tilde{f}(L)\sse \C\sm \{0\}$. If $\sigma$ has extended winding number $1$ (respectively 0), then $L$ is Hamiltonian isotopic to $L_{CL}$ (respectively $L_{Ch}$).
\end{lemma}

\begin{proof}
Suppose that $\sigma$ has extended winding number 1; the proof for the Chekanov-type case is identical. We begin by applying \Cref{lem:hamisotozerosection} so that $L$ intersects each fiber $\tilde{f}^{-1}(c)$, $c\in \sigma$, along the zero section. We want to argue that $\sigma$ can be brought to $\sigma_{Cl}$ through an area-preserving diffeomorphism, where $\sigma_{Cl}$ is the representative curve used to define $L_{Cl}$ as in \Cref{sec:std_fillings}. 

Note that by the definition of a Lagrangian filling as well as \Cref{lem:compends}, $\sigma$ and $\sigma_{Cl}$ coincide in some subset of their ends. Indeed, we have that $Ends(L)=(T_1,\infty)\times \Lambda_h$ and $Ends(L_{Cl})=(T_2,\infty)\times \Lambda_h$ for some $T_1,T_2>0$. And so choosing $T=\max \{T_1,T_2\}$, we have that
$$Ends(\sigma)\cap Ends(\sigma_{Cl}) = \tilde{f}\big(Ends(L)\big)\cap\tilde{f}\big(Ends(L_{Cl})\big) = \tilde{f}\big((T,\infty)\times \Lambda_h\big).$$

For simplicity, we can redefine $Ends(\sigma)=Ends(\sigma_{Cl}) = \tilde{f}\big((T,\infty)\times \Lambda_h\big)$, which is the union of two open line segments in $\C$.

\begin{center}
	\begin{figure}[h!]
		\centering
		
		\caption{(Left) A curve $\sigma\sse \C$ in the base of the standard Lefschetz fibration with the yellow and orange regions being the areas needed to be displaced to bring $\sigma$ to $\sigma_{Cl}$ by an area-preserving diffeomorphism. (Right) The curve $\sigma$ following the Hamiltonian isotopy that adds a bump, which results in the yellow and orange regions having equal area.\vspace{.15cm}}
		\label{fig:horiz_iso_1}
	\end{figure}
\end{center}

Since our fillings are open submanifolds in the open symplectic manifold $(B^4_1,\omega_0)$, the following idea is possible to define a Hamiltonian isotopy: adding a bump to $\sigma$ is a Hamiltonian isotopy that adds or subtracts the necessary area in order to deform $\sigma$ to $\sigma_{Cl}$ through an area-preserving diffeomorphism. Rigorously, fix coordinates $(x,y)\in \C$ where $x+iy=z$ for any $z\in\C$ and let $\beta:\R\to \R$ be an appropriate compactly supported bump function whose graph looks like the bump added in \Cref{fig:horiz_iso_1}. Then the family of Hamiltonian functions
\begin{align*}
H_t:\C\times [0,1]&\to \R\\
(x,y,t)&\mapsto t\int_{-\infty}^x \beta(\tau)d\tau,
\end{align*}
generates a compactly supported Hamiltonian flow $\phi_t(x,y) = (x,y+t\beta(x))$ that brings the $x$-axis to the graph of $\beta(x)$. Thus, with appropriate scaling and translation applied to $H_t$, we can apply this Hamiltonian isotopy to one of the open line segments of $Ends(\sigma)$ in order to add (or subtract) the correct area so that $\sigma$ can be deformed to $\sigma_{Cl}$ through an area-preserving diffeomorphism.

Given that both $\sigma$ and $\sigma_{Cl}$ have extended winding number $0$, it follows that such a diffeomorphism extends to an isotopy $\sigma_t$ such that $\sigma_0=\sigma$, $\sigma_1=\sigma_{Cl}$, and $\sigma_t$ is disjoint from $0\in \C$ for all $t$. In particular, $\sigma_t$ has extended winding number $0$ for all $t$. Therefore, the exact Lagrangian isotopy given by 
$$L_t = \tilde{f}^{-1}(\sigma_t)\cap \{||\wt{z_1}||^2-||\wt{z_2}||^2=0\}$$ 
has $L_0=L$, $L_1=L_{Cl}$, and $L_t$ is embedded for all $t$ as we avoid the unique singular fiber $\tilde{f}^{-1}(0)$. Thus we conclude that $L$ is Hamiltonian isotopic to $L_{Cl}$ through a compactly supported isotopy that is standard outside of the subset $M\sse B^4_1$.
\end{proof}

\end{document}
